\section{Limitations}
\label{sec:limitations}

The method assumes that a test image contains enough reliable normal evidence
to estimate an image-conditioned normal reference. If the whole image is
shifted, if the anomaly covers most of the object, or if normal regions are
visually ambiguous, the estimated reference can be contaminated. The current
conservative update reduces this risk but does not remove it.

The boundary-aware reliability signal mainly addresses the confusion between
high-frequency local variation and low-frequency structure boundaries. It may
be less informative for logical anomalies, global geometric anomalies, or
defects whose evidence is primarily semantic rather than local. Finally, the
mechanism conclusions should be kept within the controlled MVTec/VisA
experiments reported here. Broader dataset claims require the same controlled
ablation protocol rather than relying on system-level or externally reported
baseline tables.
